% Figure: the four rigid and scaling transformations of a base curve.
% A single base parabola y = x^2 is shifted, reflected, and scaled to
% show how the four operations move the graph. The point is that the
% engineer reads a transformed curve by undoing the transformations
% in the right order.
\begin{figure}[htbp]
\centering
\begin{tikzpicture}[scale=0.95,
  ax/.style={->, draw=engblack, thick},
  tick/.style={draw=engblack!60, thin},
  base/.style={draw=engblack, thick},
  shiftc/.style={draw=engblue, thick},
  scalec/.style={draw=engred, thick, dashed},
  refl/.style={draw=engamber, thick, dash dot},
  ann/.style={font=\footnotesize, align=left},
]
% --- axes centred ---
\draw[ax] (-3.3, 0) -- (3.5, 0) node[below right, ann] {$x$};
\draw[ax] (0, -2.3) -- (0, 4.3) node[left, ann] {$y$};
\foreach \x in {-3, -2, -1, 1, 2, 3} {
  \draw[tick] (\x, -0.05) -- (\x, 0.05);
  \node[ann, below, yshift=-2pt] at (\x, 0) {$\x$};
}
\foreach \y in {-2, -1, 1, 2, 3, 4} {
  \draw[tick] (-0.05, \y) -- (0.05, \y);
  \node[ann, left, xshift=-2pt] at (0, \y) {$\y$};
}
% --- base: y = x^2 ---
\draw[base, domain=-2:2, samples=40, smooth] plot (\x, {\x*\x});
\node[ann, anchor=south west, draw=engblack, fill=white, inner sep=2pt]
  at (1.4, 3.6) {$y = x^{2}$};
% --- horizontal shift: y = (x+1.5)^2 ---
\draw[shiftc, domain=-3.3:0.5, samples=40, smooth]
  plot (\x, {(\x+1.5)*(\x+1.5)});
\node[ann, anchor=east, draw=engblue, fill=white, inner sep=2pt]
  at (-2.0, 3.0) {$y = (x+1.5)^{2}$};
% --- vertical scale + shift down: y = 0.5 x^2 - 1.5 ---
\draw[scalec, domain=-2.8:2.8, samples=40, smooth]
  plot (\x, {0.5*\x*\x - 1.5});
\node[ann, anchor=west, draw=engred, fill=white, inner sep=2pt]
  at (1.6, 1.4) {$y = \tfrac{1}{2}x^{2} - \tfrac{3}{2}$};
% --- reflection: y = -0.6 x^2 + 3 ---
\draw[refl, domain=-2.4:2.4, samples=40, smooth]
  plot (\x, {-0.6*\x*\x + 3});
\node[ann, anchor=south, draw=engamber, fill=white, inner sep=2pt]
  at (-2.6, 1.2) {$y = -\tfrac{3}{5}x^{2}+3$};
\end{tikzpicture}
\caption[The four transformations of a base curve]{One base parabola
$y = x^{2}$ (solid black) and three transforms of it. The blue curve
$y = (x + 1.5)^{2}$ shifts the base left by $1.5$ (the inside
operation moves the graph the opposite way to its sign). The red
dashed curve $y = \tfrac{1}{2}x^{2} - \tfrac{3}{2}$ scales the base
vertically by $\tfrac{1}{2}$ and shifts it down by $\tfrac{3}{2}$
(outside operations act in the intuitive direction). The amber
dash-dotted curve $y = -\tfrac{3}{5}x^{2} + 3$ reflects the base
across the horizontal and shifts it up by $3$. The reading rule:
inside operations on $x$ act horizontally and against their sign,
outside operations on the whole expression act vertically and with
their sign.}
\label{fig:vol02:ch02:transformations}
\end{figure}
